% ============================================================
% Section 6: Intentionality Vector (VI)
% ============================================================

\section{Intentionality Vector (VI)}
\label{sec:vi}

The Intentionality Vector (VI) monitors manifold health and applies
corrections when it detects degradation. It functions as an ``immune
system'' for the manifold, preventing dimensional collapse and
confidence inflation.

\subsection{$\phi$ Field: Manifold Health}

Manifold health is quantified by the scalar field $\phi(\calM)$,
composed of 4 indicators:

\begin{equation}
    \phi(\calM) = \sum_{i=1}^{4} w_i \cdot \phi_i
    \label{eq:phi_field}
\end{equation}

with $\sum w_i = 1$. The 4 components are:

\paragraph{$\phi_1$: Dimensional Diversity.}
Measures spread across the 5 dimensions via per-axis standard deviation:

\begin{equation}
    \phi_1 = \text{clamp}\left(\frac{\overline{\text{std}(c_k)}}{0.29}, 0, 1\right)
\end{equation}

where $\overline{\text{std}(c_k)}$ is the mean of per-axis standard
deviations across the batch. The normalizer $0.29$ corresponds to the
std of a uniform $[0,1]$ distribution. Low values indicate
\textbf{dimensional collapse}.

\paragraph{$\phi_2$: Dispersion.}
Measures distance diversity from the batch centroid:

\begin{equation}
    \phi_2 = \text{clamp}\left(\frac{\|\bm{c} - \bar{\bm{c}}\|_2}{\sqrt{5}}, 0, 1\right)
\end{equation}

where $\bar{\bm{c}}$ is the centroid and $\sqrt{5}$ is the maximum
L2 distance in $[0,1]^5$.

\paragraph{$\phi_3$: Entropy.}
Measures distributional uniformity of coordinates across axes:

\begin{equation}
    p_k = \frac{\bar{c}_k}{\sum_j \bar{c}_j + \epsilon}, \quad
    \phi_3 = \text{clamp}\left(\frac{-\sum_k p_k \log p_k}{\log 5}, 0, 1\right)
\end{equation}

\paragraph{$\phi_4$: Confidence Distribution.}
Uses a U-shaped profile with ideal standard deviation $\sigma^* = 0.15$:

\begin{equation}
    \phi_4 = \text{clamp}\left(1 - \frac{|\text{std}(\kappa) - \sigma^*|}{0.5}, 0, 1\right)
\end{equation}

This penalizes both uniformly high confidence (excess) and uniformly
low confidence (collapse). The normalizer $0.5$ provides a wider
tolerance range than $\sigma^*$.

\subsection{Severity and Activation}

Severity is computed as a 70/30 blend of analytical and learned
pathways:

\begin{equation}
    s_{\text{anal}} = \sqrt{\max\left(0, 1 - \frac{\phi}{\phi_{\text{critical}}}\right)}
    \label{eq:vi_severity_anal}
\end{equation}

\begin{equation}
    s_{\text{learn}} = \sigma\big(\text{MLP}_{4 \to 8 \to 1}(\phi_1, \phi_2, \phi_3, \phi_4)\big)
\end{equation}

\begin{equation}
    s = 0.7 \cdot s_{\text{anal}} + 0.3 \cdot s_{\text{learn}}
    \label{eq:vi_severity}
\end{equation}

with $\phi_{\text{critical}} = 0.5$ by default. The analytical pathway
provides a deterministic baseline via the square root curve, while
the learned pathway allows the model to adapt severity to specific
component patterns.

\subsection{Manifold Correction}

When $s > 0$, the VI computes a correction direction via a learned
MLP:

\begin{equation}
    \bm{d}_{\text{raw}} = \tanh\big(\text{MLP}_{9 \to 10 \to 5}([\phi_1, \phi_2, \phi_3, \phi_4, \bm{c}])\big)
\end{equation}

\begin{equation}
    \bm{v} = \frac{\bm{d}_{\text{raw}}}{\|\bm{d}_{\text{raw}}\| + \epsilon} \cdot s \cdot \alpha
    \label{eq:vi_correction}
\end{equation}

where $\alpha = 0.6$ is the injection strength. The direction network
takes the 4 $\phi$ components and the current 5D coordinates as input
(9 dimensions total), allowing context-dependent correction rather
than a fixed target.

\subsection{Confidence Penalty}

When confidence is inflated (high confidence with degraded manifold),
the VI applies a penalty:

\begin{equation}
    \kappa_{\text{corrected}} = \kappa \cdot (1 - \beta \cdot s)
    \label{eq:vi_conf_penalty}
\end{equation}

with $\beta = 0.4$. This prevents the model from reporting high
confidence when the manifold is in a state of collapse.

\subsection{VI Regularization}

The VI regularization loss encourages minimum health and penalizes
excessive severity:

\begin{equation}
    \calL_{\text{VI}} = \E\big[\text{ReLU}(\phi_{\text{critical}} - \phi)^2\big]
    + 0.3 \cdot \E[s^2]
\end{equation}

with $\phi_{\text{critical}} = 0.5$. The first term ensures the manifold
does not degenerate during training. The second term penalizes high
severity, encouraging the model to maintain health proactively rather
than relying on corrections.

\subsection{Preliminary Evidence: Collapse Prevention}

In simulations with the ATIC system (deterministic predecessor, where VI
operates as a post-hoc correction module), the VI demonstrated:

\begin{itemize}
    \item Effective dimensionality recovery: $D_{\text{eff}}$
    from 2.2 to 3.2 (out of 5 possible)
    \item Inflated confidence reduction: from 0.60 to 0.57
    \item 3 out of 4 health indicators improved
\end{itemize}

\textbf{Note:} These results are from the deterministic ATIC system, not
from \aletheion{} training. In \aletheion{}, these mechanisms are trained
end-to-end, which may yield different dynamics. Empirical validation on
the trained 354M model is ongoing (see \Cref{sec:conclusion}, Future Work).
